\chapter*{Abstract}
\addcontentsline{toc}{chapter}{Abstract}

Students now routinely use generative AI systems even for complex tasks, which carries the risk that
solutions are adopted without scrutiny. This thesis investigates whether an AI system specifically
designed for reflection (a ``Reflexionsbot'') leads students to work through process decisions in
software process management more independently and more reflectively than a generic AI system does.

To this end, a Reflexionsbot was implemented that guides students through Socratic questioning and
adaptive response depth instead of direct answers. In a mini-experiment conducted within the course
``Grundlagen des Software Engineering 1'' in the Software Engineering bachelor's programme at
Heilbronn University, one group worked through two realistic process scenarios using the
Reflexionsbot, while a comparison group used a generic AI system. The resulting 31 submissions and
the corresponding chat logs were coded using structuring qualitative content analysis based on Hatton
and Smith's model of reflective levels, extended by text authorship as a separate category, and
triangulated with a pre/post survey (n=22 and n=20, respectively).

The share of submissions without clearly student authorship is lower in the Reflexionsbot group
(8 of 16 versus 12 of 15). This gap is partly induced by the task requirement imposed on the
comparison group and is considerably smaller when only clearly AI-generated cases are considered
(7 of 16 versus 8 of 15). The direct learning measures do not support an advantage of the
Reflexionsbot. A ceiling effect in the knowledge test precludes any conclusion, and the demonstrated
depth of reflection is descriptively even lower than in the comparison group (mean level 0.75 versus
1.33), although with three versus eight codable cases this difference carries no directional weight.
Only students' self-assessment of reasoning and weighing is more favourable. Didactic design
therefore influences usage behaviour, but a learning advantage does not follow automatically. What
matters is students' willingness to engage with the dialogue.

\vspace{1em}
\noindent\textbf{Keywords:} Generative AI, Reflection Bot, Socratic Questioning, Software Engineering
Education, Mixed Methods, Qualitative Content Analysis

